\documentclass[a4paper,13pt,3p,twoside]{report}
\usepackage{scrextend}
\changefontsizes{13pt}
\usepackage[utf8]{vietnam}
\usepackage[top=2cm, bottom=2cm, left=3.5cm, right=2.5cm]{geometry}

\usepackage{graphicx} % Cho phép chèn hình ảnh
\usepackage{fancybox} % Tạo khung box
\usepackage{indentfirst} % Thụt đầu dòng ở dòng đầu tiên trong đoạn
\usepackage{amsthm} % Cho phép thêm các môi trường định nghĩa
\usepackage{latexsym} % Các kí hiệu toán học
\usepackage{amsmath} % Hỗ trợ một số biểu thức toán học
\usepackage{amssymb} % Bổ sung thêm kí hiệu về toán học
\usepackage{amsbsy} % Hỗ trợ các kí hiệu in đậm
\usepackage{times} % Chọn font Times New Roman
\usepackage{array} % Tạo bảng array
\usepackage{enumitem} % Cho phép thay đổi kí hiệu của list
\usepackage{subfiles} % Chèn các file nhỏ, mỗi chương một file
\usepackage{titlesec} % Chỉnh sửa tiêu đề chương, mục
\usepackage{titletoc}
\usepackage{chngcntr} % Thiết lập lại cách đánh số caption
\usepackage{pdflscape} % Bảng đặt theo chiều ngang giấy
\usepackage{afterpage}
\usepackage[ruled,vlined]{algorithm2e} % Viết giải thuật
\usepackage{capt-of} % Caption lớn trên trang ngang
\usepackage{multirow} % Gộp ô
\usepackage{longtable} % Bảng dài nhiều trang, lặp header
\usepackage{fancyhdr} % Tùy biến header và footer
\usepackage{appendix}

\usepackage[font=small,labelfont=bf,justification=centering]{caption}

\usepackage{listings}
\usepackage{float}
\usepackage{subcaption}
\usepackage{xurl}
\usepackage{setspace}
\usepackage{parskip}
\usepackage{tocbasic}

% ===================================================

\usepackage[backend=biber,style=ieee]{biblatex}
\addbibresource{Danh_sach_tai_lieu_tham_khao.bib}
\setcounter{biburlnumpenalty}{9000} % Cho phép xuống dòng trong URL
\setcounter{biburlucpenalty}{9000}
\setcounter{biburllcpenalty}{9000}

\renewcommand\appendixname{PHỤ LỤC}
\renewcommand\appendixpagename{PHỤ LỤC}
\renewcommand\appendixtocname{PHỤ LỤC}

% ===================================================

\fancypagestyle{plain}{%
\fancyhf{} % Xóa toàn bộ header và footer
\fancyfoot[RO,RE]{\thepage}
\renewcommand{\headrulewidth}{0pt}
\renewcommand{\footrulewidth}{0pt}}

\setlength{\headheight}{10pt}

\def \TITLE{ĐỒ ÁN TỐT NGHIỆP}
\def \AUTHOR{ĐỖ VĂN VIỆT}
\newcommand{\TenDeTai}{Hệ thống chấm công nhận diện khuôn mặt ứng dụng TinyML trên vi điều khiển ESP32-S3}
\newcommand{\GVHD}{TS. Đỗ Công Thuần}
\newcommand{\ThoiGianNop}{HÀ NỘI, 01/2027}

% ===================================================
\titleformat
    {\chapter} % command
    [hang] % shape
    {\centering\bfseries} % format
    {CHƯƠNG \thechapter.\ } % label
    {0pt} % sep
    {} % before
    [] % after
\titlespacing*{\chapter}{0pt}{-20pt}{20pt}

\titleformat
    {\section} % command
    [hang] % shape
    {\bfseries} % format
    {\thechapter.\arabic{section}\ \ \ \ } % label
    {0pt} % sep
    {} % before
    [] % after
\titlespacing{\section}{0pt}{\parskip}{0.5\parskip}

\titleformat
    {\subsection} % command
    [hang] % shape
    {\bfseries} % format
    {\thechapter.\arabic{section}.\arabic{subsection}\ \ \ \ } % label
    {0pt} % sep
    {} % before
    [] % after
\titlespacing{\subsection}{30pt}{\parskip}{0.5\parskip}

\renewcommand\thesubsubsection{\alph{subsubsection}}
\titleformat
    {\subsubsection} % command
    [hang] % shape
    {\bfseries} % format
    {\alph{subsubsection}, \ } % label
    {0pt} % sep
    {} % before
    [] % after
\titlespacing{\subsubsection}{50pt}{\parskip}{0.5\parskip}

% ===================================================
\usepackage{hyperref}
\hypersetup{pdfborder = {0 0 0}}
\hypersetup{pdftitle={\TITLE},
	pdfauthor={\AUTHOR}}

\usepackage[all]{hypcap} % Tham chiếu chính xác đến hình và bảng

\graphicspath{{Hinhve/}{../Hinhve/}} % Thư mục chứa hình

\counterwithin{figure}{chapter} % Đánh số hình theo chương: Hình 1.1, 1.2

\title{\bf \TITLE}
\author{\AUTHOR}

\setcounter{secnumdepth}{3} % Cho phép subsubsection trong report

\theoremstyle{definition}
\newtheorem{example}{Ví dụ}[chapter]

\onehalfspacing
\setlength{\parskip}{6pt} % Khoảng cách giữa hai đoạn
\setlength{\parindent}{15pt} % Lùi đầu dòng

% Khoảng cách hình, bảng với văn bản cố định và đồng nhất mọi nơi.
\raggedbottom
\setlength{\intextsep}{8pt}
\setlength{\textfloatsep}{8pt}
\setlength{\floatsep}{8pt}
\setlength{\abovecaptionskip}{5pt}
\setlength{\belowcaptionskip}{2pt}

% =========================== BODY ===============
\begin{document}
\subfile{Bia}
\subfile{Bia_Lot}

% ===================================================
\newpage
\pagenumbering{gobble}
\subfile{CHUONG/01_Loi_cam_on}

\newpage
\pagenumbering{gobble}
\subfile{CHUONG/02_Tom_tat_noi_dung_tieng_Viet}

% ===================================================
\newpage
\pagenumbering{roman}
\renewcommand*\contentsname{MỤC LỤC}

\titlecontents{chapter}
    [0.0cm] % left margin
    {\bfseries\vspace{0.3cm}} % above code
    {{\bfseries{\scshape}
    CHƯƠNG \thecontentslabel.\ }} % numbered format
    {} % unnumbered format
    {\titlerule*[0.3pc]{.}\contentspage} % filler-page-format

\titlecontents{section}
    [0.0cm] % left margin
    {\vspace{0.3cm}} % above code
    {\thecontentslabel \ } % numbered format
    {} % unnumbered format
    {\titlerule*[0.3pc]{.}\contentspage} % filler-page-format

\titlecontents{subsection}
    [1.0cm] % left margin
    {\vspace{0.3cm}} % above code
    {\thecontentslabel \ } % numbered format
    {} % unnumbered format
    {\titlerule*[0.3pc]{.}\contentspage} % filler-page-format

\addtocontents{toc}{\protect\thispagestyle{empty}}
\tableofcontents
\thispagestyle{empty}
\cleardoublepage

\renewcommand{\listfigurename}{DANH MỤC HÌNH VẼ}
{\let\oldnumberline\numberline
\renewcommand{\numberline}{Hình~\oldnumberline}
\listoffigures}
\newpage

\renewcommand{\listtablename}{DANH MỤC BẢNG BIỂU}
{\let\oldnumberline\numberline
\renewcommand{\numberline}{Bảng~\oldnumberline}
\listoftables}

\newpage
\subfile{CHUONG/03_Danh_muc_viet_tat}

\newpage
\subfile{CHUONG/04_Thuat_ngu}
% ===================================================

\newpage
\pagenumbering{arabic}

\pagestyle{fancy}
\fancyhf{}
\fancyhead[RE, LO]{\leftmark}
\fancyfoot[RE, LO]{\thepage}

\chapter{GIỚI THIỆU ĐỀ TÀI}
\label{chapter:Introduction}
\subfile{CHUONG/1_Gioi_Thieu_De_Tai}

\newpage
\chapter{CƠ SỞ LÝ THUYẾT}
\label{chapter:Related_works}
\subfile{CHUONG/2_Co_so_ly_thuyet}

\newpage
\chapter{PHÂN TÍCH THIẾT KẾ HỆ THỐNG}
\label{chapter:Design}
\subfile{CHUONG/3_Thiet_ke}

\newpage
\chapter{TRIỂN KHAI HỆ THỐNG}
\label{chapter:Implementation}
\subfile{CHUONG/4_Trien_khai}

\newpage
\chapter{THỰC NGHIỆM VÀ ĐÁNH GIÁ}
\label{chapter:Experiment}
\subfile{CHUONG/5_Thuc_nghiem_va_danh_gia}

\newpage
\chapter*{KẾT LUẬN}
\addcontentsline{toc}{chapter}{KẾT LUẬN}
\markboth{KẾT LUẬN}{}
\subfile{CHUONG/6_Ket_luan}

% ===================================================
\newpage
\printbibliography[heading=bibintoc, title={TÀI LIỆU THAM KHẢO}]

\appendixpage
\appendices
\addappheadtotoc

\titleformat{\chapter}[hang]{\centering\bfseries}{ \thechapter.\ }{0pt}{}[]
\titlespacing*{\chapter}{0pt}{-20pt}{20pt}

\titlecontents{chapter}
    [0.0cm] % left margin
    {\bfseries\vspace{0.3cm}} % above code
    {{\bfseries{\scshape} \thecontentslabel.\ }} % numbered format
    {} % unnumbered format
    {\titlerule*[0.3pc]{.}\contentspage} % filler-page-format
\subfile{CHUONG/phu_luc}
\end{document}
